\documentclass[11pt]{article}
\usepackage[margin=25mm]{geometry}
\usepackage{amsmath,amssymb}
\newcommand{\AtelierReals}{\ensuremath{\mathbb{R}}}
\newcommand{\AtelierPartial}[2]{\frac{\partial #1}{\partial #2}}
\newcommand{\AtelierNote}[2][Note]{\textbf{#1:} #2}
\NewDocumentCommand{\AtelierNorm}{s m}{%
 \IfBooleanTF{#1}{\left\lVert #2\right\rVert}{\lVert #2\rVert}}
\ExplSyntaxOn
\cs_new_protected:Npn \atelier_items:n #1
 {
  \begin{itemize}
  \clist_map_inline:nn {#1} {\item ##1}
  \end{itemize}
 }
\NewDocumentCommand{\AtelierItems}{m}{\atelier_items:n {#1}}
\ExplSyntaxOff
\begin{document}
\title{A guided macro example}
\author{}
\date{}
\maketitle
\section{Give notation a name}
The set of real numbers is \AtelierReals.
Inside mathematics the same command works: $x\in\AtelierReals$.

Let $U(c,s)=\log c+s$ with $c>0$. The first and second arguments
of the partial-derivative macro are the function and its variable:
\[
\AtelierPartial{U}{c}=\frac{1}{c},
\qquad
\AtelierPartial{U}{s}=1.
\]

\section{Use an optional label}
\AtelierNote{The first argument has a default value.}

\AtelierNote[Input]{An explicit label replaces that default.}

\section{Read a star}
Let $x$ and $y$ be vectors in the same space.
The unstarred norm has ordinary delimiters:
\[
\AtelierNorm{x+y}.
\]
The starred call sizes delimiters to the taller argument:
\[
\AtelierNorm*{\frac{x+y}{2}}.
\]
The star state is argument one; the norm's content is argument two.

\section{Map a comma list}
The public command below passes a comma list to an internal function:
\AtelierItems{Input,Calculation,Output}
Inside the enclosing function definition, the current list item
uses a doubled parameter marker. Removing one input word should
remove one list item without changing the function.
\end{document}
